
\documentclass[%
pdftex,
oneside,			% Einseitiger Druck.
11pt,				% Schriftgroesse
parskip=half,		% Halbe Zeile Abstand zwischen Absätzen.
%	topmargin = 10pt,	% Abstand Seitenrand (Std:1in) zu Kopfzeile [laut log: unused]
headheight = 12pt,	% Höhe der Kopfzeile
%	headsep = 30pt,	% Abstand zwischen Kopfzeile und Text Body  [laut log: unused]
headsepline,		% Linie nach Kopfzeile.
footsepline,		% Linie vor Fusszeile.
footheight = 16pt,	% Höhe der Fusszeile
abstracton,		% Abstract Überschriften
DIV=calc,		% Satzspiegel berechnen
BCOR=8mm,		% Bindekorrektur links: 8mm
headinclude=false,	% Kopfzeile nicht in den Satzspiegel einbeziehen
footinclude=false,	% Fußzeile nicht in den Satzspiegel einbeziehen
listof=totoc,		% Abbildungs-/ Tabellenverzeichnis im Inhaltsverzeichnis darstellen
toc=bibliography,	% Literaturverzeichnis im Inhaltsverzeichnis darstellen
]{scrreprt}	% Koma-Script report-Klasse, fuer laengere Bachelorarbeiten alternativ auch: scrbook


\def \titel {Use of generative AI systems to support processes, workflows, or organizations} 
\def \arbeit {Project Report}
\def \autor {Jasper Bremenkamp}
\def \abschluss {M.Sc.}
\def \studiengang {Data Science}
\def \datumAbgabe {\today}
\def \modul {Project: AI Excellence with Creative Prompt Techniques}
\def \matrikelnr {92125193}
\def \tutor {Marek Malucha}

\def \seitenrand {2cm}

% PDF Einstellungen
\usepackage{hyperref}

\hypersetup{ 
    unicode = true,
    pdftitle={\titel},
    pdfauthor={\autor},
    pdfsubject={\arbeit},
    pdfcreator={pdflatex, LaTeX with KOMA-Script},
    pdfpagemode=UseOutlines,    % Beim Oeffnen Inhaltsverzeichnis anzeigen
    pdfdisplaydoctitle=true,    % Dokumenttitel statt Dateiname anzeigen.
    pdflang={en},   % Sprache des Dokuments.
    hidelinks
}

\usepackage[left=2cm,right=2cm,top=2cm,bottom=2cm]{geometry} % margins	% Seitenränder und Abstände
\usepackage[activate]{microtype} %Zeilenumbruch und mehr
\usepackage[onehalfspacing]{setspace}
\usepackage{makeidx}
\usepackage[autostyle=true,german=quotes]{csquotes}
\MakeOuterQuote{"} % Für deutsche Anführungszeichen 
\usepackage{longtable}
\usepackage{enumitem}	% mehr Optionen bei Aufzählungen
\usepackage{graphicx}
\usepackage{pdfpages}   % zum Einbinden von PDFs
\usepackage{xcolor} 	% für HTML-Notation
\usepackage{float}
\usepackage{array}
\usepackage{calc}		% zum Rechnen (Bildtabelle in Deckblatt)
\usepackage[right]{eurosym}
\usepackage{wrapfig}
\usepackage{pgffor} % für automatische Kapiteldateieinbindung
\usepackage[perpage, hang, multiple, stable]{footmisc} % Fussnoten
\usepackage[printonlyused]{acronym} % falls gewünscht kann die Option footnote eingefügt werden, dann wird die Erklärung nicht inline sondern in einer Fußnote dargestellt
\usepackage{listings}
\usepackage{bookmark}
\usepackage[english]{babel}
\usepackage[utf8]{inputenc}
\usepackage{scrhack}
\usepackage[T1]{fontenc} %Damit Umlaute auch beim Kopieren richtig dargestellt werden
\usepackage{caption} %für captionof
\usepackage{graphicx}
\usepackage{placeins}


    %%%%%%% Neues Design %%%%%%%%%%%

    \usepackage{helvet} %Arial-Ersatz
    \renewcommand{\familydefault}{\sfdefault}
    
    \def\UrlFont{\rm} % Print URLs not in Typewriter Font

    \usepackage{etoolbox}
    \makeatletter
    \patchcmd{\scr@startchapter}{\if@openright\cleardoublepage\else\clearpage\fi}{}{}{} %entfernt das starten einer neuen Seite bei einem neuen Chapter
    \makeatother

    %\addtolength{\footskip}{-0.7cm}% foot larger by 0,7 cm  (Raises the page number)
    %\setlength{\parindent}{6pt} % Indent at start of paragraphs  6pt

    \usepackage{titlesec} % define size for section headings
    \titleformat{\chapter}{\normalfont\fontsize{12pt}{1.5}\bfseries}{\thechapter}{1em}{}
    \titleformat{\section}{\normalfont\fontsize{12pt}{1.5}\bfseries}{\thesection}{1em}{}
    \titleformat{\subsection}{\normalfont\fontsize{12pt}{1.5}\bfseries}{\thesubsection}{1em}{} 

    \renewcommand{\footnotesize}{\fontsize{10pt}{1.5pt}\selectfont}

    \titlespacing{\chapter}{0pt}{10pt}{10pt}

    %%%%%%%%%%%%%%%%%%%%%%%%%%%%%%%

% Notizen. Einsatz mit \todo{Notiz} oder \todo[inline]{Notiz}. 
\usepackage[obeyFinal,backgroundcolor=yellow,linecolor=black]{todonotes}
% Alle Notizen ausblenden mit der Option "final" in \documentclass[...] oder durch das auskommentieren folgender Zeile
% \usepackage[disable]{todonotes}

% Kommentarumgebung. Einsatz mit \comment{}. Alle Kommentare ausblenden mit dem Auskommentieren der folgenden und dem aktivieren der nächsten Zeile.
\newcommand{\comment}[1]{\par {\bfseries \color{blue} #1 \par}} %Kommentar anzeigen

% der neue Zitierstil ist Käse
\global\def\VAlterZitierstil{true}

\global\def\VTrueValue{true}

\ifx\VAlterZitierstil\VTrueValue
    % \usepackage[
    %     backend=biber,     % biber muss seperat installiert werden
    %     bibwarn=true,
    %     bibencoding=utf8,   % wenn .bib in utf8, sonst ascii
    %     sortlocale=de_DE,
    %     sorting=none,
    %     style=apa %apa instead of IEEE
    % ]{biblatex}

    \usepackage[backend=biber, style=apa, pagetracker, ibidtracker=constrict, apamaxprtauth=20 ]{biblatex}

    %Enable IBAN on APA style
    \usepackage{xpatch}
    \xpretobibmacro{doi+url}{\printfield{isbn}}{}{}

    \addbibresource{bibliographie.bib}
\else
    %eigentlich totaler Käse, kann nach Bewertung der Arbeit für zukünftige Arbeiten entfernt werden

    % New IUBH Cite Style / Old Cite style over this block
    % \usepackage{natbib}
    % \setcitestyle{aysep={}} % remove comma as delimiter 

    \usepackage{apacite}

    %search and replace von apacite zur biblatex
    %Search: \\citeNP<([^}]*)>\[([^}]*)\]\{([^}]+)\}
    %Replace: \\cite[$1][$2]{$3}
\fi

\definecolor{codegreen}{rgb}{0,0.6,0}
\definecolor{codegray}{rgb}{0.5,0.5,0.5}
\definecolor{codepurple}{rgb}{0.58,0,0.82}
\definecolor{backcolour}{rgb}{0.95,0.95,0.92}

\lstdefinestyle{mystyle}{
    backgroundcolor=\color{backcolour},
    commentstyle=\color{codegreen},
    keywordstyle=\color{magenta},
    numberstyle=\tiny\color{codegray},
    stringstyle=\color{codepurple},
    basicstyle=\ttfamily\footnotesize,
    breakatwhitespace=false,
    breaklines=true, 
    captionpos=b,
    keepspaces=true,
    numbers=left,
    numbersep=5pt,
    showspaces=false,
    showstringspaces=false,
    showtabs=false,
    tabsize=2
}

\lstdefinestyle{yaml}{
     basicstyle=\color{blue}\footnotesize,
     rulecolor=\color{black},
     string=[s]{'}{'},
     stringstyle=\color{blue},
     comment=[l]{:},
     commentstyle=\color{black},
     morecomment=[l]{-}
 }

\lstset{style=mystyle}

\renewcommand{\lstlistingname}{Code snippets}% 
\renewcommand{\lstlistlistingname}{List of code snippets}%

\begin{document}

\pagenumbering{Roman}

\begin{titlepage}
    \begin{figure}[h]
        \centering
        \advance\leftskip--2cm
        \includegraphics[width=0.50\textwidth]{assets/logo.pdf}
        \end{figure}
        \FloatBarrier
        
        \begin{Large} 
        \begin{center}
        \arbeit
        \end{center}
        \end{Large} 
        
        \vspace*{30mm}
        
        \begin{Large} 
        \begin{center}
        \textbf{\titel}
        \end{center}
        \end{Large}
        
        \vspace*{45mm}
        
        \begin{large} 
        \begin{center}
        \autor
        \end{center}
        \end{large} 
        
        \vspace*{-6mm}
        
        \begin{large} 
        \begin{center}
        Matriculation number: \matrikelnr
        \end{center}
        \end{large} 
        
        \vspace*{-6mm}
        
        \begin{large} 
        \begin{center}
        % ursprünglich Straße
        \end{center}
        \end{large} 
        
        \vspace*{24mm}
        
        \begin{large} 
        \begin{center}
        is done in the course \enquote{\modul} \\
        of the degree program \abschluss { }\studiengang
        \end{center}
        \end{large} 
        
        \vspace*{10mm}
        
        \begin{large} 
        \begin{center}
        Advisor: \tutor 
        \end{center}
        \end{large}         
        \vspace*{-6mm}
        
        \begin{large} 
        \begin{center}
        Delivery date: \datumAbgabe
        \end{center}
        \end{large} 
\end{titlepage}

\renewcommand{\abstractname}{Abstract}
\begin{abstract}


    %Evtl. noch was zum Ergebnis schreiben: ein oder zwei Sätze


    %Teil des Ergebnisses muss hier noch stehen


    %Draft! Durchgehen und  überarbeiten!
    %%%%%%%%%%%%%%%%%%%%

    To-Do
%%%%%


\end{abstract}

% Inhaltsverzeichnis
\begin{spacing}{1.1}
    \begingroup
    
        % auskommentieren für Seitenzahlen unter Inhaltsverzeichnis
        \renewcommand*{\chapterpagestyle}{empty}
        \pagestyle{empty}

        \setcounter{tocdepth}{3}
        %für die Anzeige von Unterkapiteln im Inhaltsverzeichnis
        %\setcounter{tocdepth}{2}
        
        \tableofcontents
        \clearpage
    \endgroup
\end{spacing}
\newpage

%Abbildungsverzeichnis
\cleardoublepage 
\listoffigures
\clearpage

\lstlistoflistings %List of listings
\clearpage

%Tabellenverzeichnis
%\cleardoublepage
%\listoftables

\chapter{Introduction}
\pagenumbering{arabic}


    What was/were the goal/goals of your AI project?

    Which AI approach did you wanted to integrate?

    What preparatory actions did you take beforehand?

    Which generative AI methodologies and tools did you select and why?

    Why did you opt for this AI approach?

    % Ziel: Einleitung: ca. 20%
    % Dafür wären eigentlich 2,5 Seiten nötig! Evtl strecken oder Bilder einbauen!

    \newpage

\chapter{Implementation of the AI Project} 

    Project structure plan
    \begin{itemize}
        \item Outline and scheduling of significant dates and events you have planned
        \item  Project resources: resource/cost planning (if applicable)
    \end{itemize}

    Project progress:
    \begin{itemize}
        \item How did you implement AI methodologies and tools in your project?
        \item Which prompts did you use?
        \item How did you develop, organize, and fine-tune your prompts?
        \item How did you implement the usage of prompts in your project field?
        \item What processes did you implement in your project for the integration of AI into your project area?
        \item Actual Stakeholder, risk, or ethical analysis, if applicable
    \end{itemize}

\chapter{Achievement of objectives}

    \begin{itemize}
        \item What result were achieved?
        \item What impacts and changes were realized and measured through the introduction/ implementation of generative AI?
        \item Does the result of your AI integration/project correspond to you goal? If not, why not?
        \item What conclusions can be drawn from the integration of AI during the project?
        \item Can a sustainable structure be created to anchor the usage of generative AI?
        \item What can be improved in the implementation of AI?
        \item What changes and/or continuous improvements need to be implemented and how?
    \end{itemize}

\chapter{Conclusion}

    \begin{itemize}
        \item Reflection and future goals
        \item Possible development for the used AI approach
        \item Possible development for the used AI methodologies and tools
        \item Ethical reflectiont.
    \end{itemize}

    % Schluss ca. 10%

% 12-15 Seiten Inhalt
% Tabellen- und Abbildungsverzeichnis einfügen

% Literaturverzeichnis
\cleardoublepage

\ifx\VAlterZitierstil\VTrueValue
    \printbibliography
\else
    %\bibliographystyle{iubh}
    \bibliographystyle{apacite} %Einträge sind unsortiert
    \bibliography{bibliographie.bib}
\fi

% To-Dos:
% Titelfolie nochaml überprüfen:
% wegen Übersetzung und ob das so Okay ist auf Englisch

\end{document}
